\section{The Reservation Discipline}

The discipline adds one object and one decision to the serving stack. The
object is the reservation: a ledger entry binding a flow to a guaranteed
resource vector for a lifetime. The decision is admission: whether a flow
enters the small reserved subset, enters an aggregate service class, or is
declined early. Everything else in this section is the precise semantics
of those two things, together with the supervisory and trust elements that
complete the contract. We emphasize at the outset what the discipline
keeps: the gateway accounting of Section 2 remains the identity and
metering substrate, and the scheduling engines of Section 2 remain the
execution machinery. The discipline sits at admission and constrains what
those layers may do to an admitted flow; it replaces neither.

\subsection{Admission and the reservation vector}

A flow arrives with a descriptor: an estimated invocation count or
duration, a requested cumulative output-token budget, an estimated cache
footprint, a service class, a provenance class (Section 4.4), and
optionally a deadline. Estimates may be wrong; the semantics below are
defined so that wrong estimates degrade a flow's own outcome, never the
guarantee already granted to another.

Admission to the reserved subset requires five conditions jointly. First,
the reserved subset on the target pool holds fewer than K flows, where K
is a deliberately small bound relative to concurrent traffic: the
discipline rations per-flow state to the flows whose requirements justify
it, and class-serves the rest, which is the surviving shape of Section
3's first lesson. Second, the pool's projected commitment, the sum over
active and admitted-but-idle reservations of their remaining resource
claims over a forecast window, remains under a configured ceiling after
adding this flow. Third, a device set exists that can hold the flow's
cache footprint under the pinning rules below; this is an affinity
feasibility test, not a best-effort placement. Fourth, the requested
token budget fits within the flow's class quota. Fifth, the flow's
provenance class is reservation-eligible (Section 4.4). A flow failing
any condition is admitted to an aggregate class, or receives an early
decline, optionally with a time quote for when admission would succeed.
Declining at the door is a feature: it is the cheapest moment in a flow's
life to say no.

Projecting compute commitment for autoregressive work is estimation, not
measurement: output lengths are unknown at admission, and the forecast
inherits that uncertainty. The discipline absorbs estimation error
structurally rather than pretending precision: the commitment ceiling is
set with margin, misestimation is charged to the misestimating flow's own
disposition (Section 10), and the forecast window's construction is an
evaluation variable, not a claimed contribution.

An admitted reservation jointly covers three resources, and the jointness
is the point: a flow starved on any one of them fails entirely, so a
guarantee over a subset is not a guarantee. The first is compute: a share
sufficient to meet the flow's class latency budget at each invocation
start. The second is a cumulative output-token budget metered at decode
across every invocation of the flow. The budget has a dual character that
single-purpose mechanisms lack: it is a floor, in that budgeted capacity
is committed and the flow cannot be silently starved out of it, and
simultaneously a bound, in that the flow cannot exceed it, which makes
the runaway loop of Section 2 a contained event rather than a discovered
one. The third is key-value cache: the flow's accumulated cache is
pinned, exempt from eviction and migration for the reservation lifetime;
pinning applies to the flow's private suffix blocks, while prefix blocks
shared across flows under paged cache management remain governed by
reference counting, which never reclaims a block a pinned flow still
references. The pin holds across inter-invocation idle gaps up to a
per-reservation idle bound.

\subsection{The guarantee, its cost surfaces, and ordinary traffic}

The guarantee is stated negatively, because that is what makes it
checkable: for the reservation lifetime, no resource-motivated action of
the serving stack may preempt the flow's invocations, evict or migrate
its pinned cache, or reclaim its committed budget. The sole override is a
halt issued by the management plane (Section 4.3), which is a safety
revocation under a designated authority, and definitionally not scheduler
preemption: no capacity condition can trigger it. This distinction is
load-bearing. A guarantee with a capacity-triggered escape clause is a
priority, and Section 2's schedulers already have priorities; the
discipline exists because priorities are not contracts.

Pinning across idle gaps is where the guarantee meets its cost, and we
give two embodiments. Under strict pinning the cache stays resident in
accelerator memory throughout, idle gaps included: maximal simplicity,
maximal stranding, appropriate when gaps are short or the reserved subset
is tiny. Under tiered pinning the cache is resident during active
invocations, and between invocations it may be moved to a guaranteed
spill tier, host memory with a reserved restore path, under a contractual
bound on restore latency and with restore capacity itself reserved. The
guarantee survives the movement: the flow observes at most the restore
bound as added start latency, and nothing it has accumulated is ever
lost. Tiered pinning is the engineering answer to the idle-gap objection
of Section 2, and the evaluation of Sections 5 through 7 prices both
embodiments without favor.

Budget exhaustion is a defined event, not an ambush. Policies are
configured per class from a small vocabulary: warn at a threshold
fraction and continue; hard-stop generation at the bound; escalate to the
management plane for a disposition; or renegotiate, which is a fresh
admission test at current headroom for an extended budget, with no
privilege carried over from the original grant. The contrast with the
gateway ceiling of Section 2 is the timing and the contract: a flow under
this discipline knows its bound at admission, progresses under a
guarantee toward it, and reaches exhaustion as a planned disposition,
rather than spending freely until a rejection converts the entire prior
spend into waste.

All other traffic runs in aggregate service classes: a small set, each
with a latency budget and a scheduling weight, imposing no per-flow state
on the serving core. Classes are where the discipline meets ordinary
traffic, and deliberately nothing here is novel: weighted scheduling
against class latency budgets is Section 2's scheduling family doing what
it does well.

Degradation is the discipline's pressure-release surface, and its signal
is the departure from precedent flagged in Section 3's second lesson.
Define committed headroom H on a pool as one minus the ratio of forward
commitments, from the reservation ledger over a forecast window, to
capacity. As H shrinks, best-effort traffic experiences actions drawn
from a configured ladder, each applied with probability p(H), monotone
increasing as H falls, beginning well before exhaustion: early admission
decline, model-tier downgrade, context truncation, token-budget haircut,
cache spill, reduced generation effort, including runtime quality
adaptation of the model itself \citep{morphserve2025}. The curve is
class-differentiated, lower classes degrading earlier and harder, and
Section 4.4 adds provenance to the keying. The signal choice is a design
claim the evaluation tests directly: queue depth reports congestion that
has already arrived; the ledger reports congestion that has been
promised; whether the forward signal is worth anything is answered,
against this design's expectation, in Sections 7 and 9.

Reservations are leases. Each carries a lifetime; renewal is an admission
test, not an entitlement; expiry returns the resources to the pool with a
grace disposition for the flow. In multi-tenant deployments the structure
nests: K, commitment ceilings, and class quotas are set per tenant
against tenant-level allocations, and the account-level provisioned
throughput of Section 2 composes naturally as the tenant's envelope
within which flow-level contracts are written. Appendix A states the
admission test and the ledger invariants at pseudocode level; the
invariants, no over-commitment of compute, memory, or budget under any
interleaving of admissions, renewals, expiries, and halts, are the
correctness core of the discipline, and the evaluation's engine enforces
them as hard test failures (Section 6).

\subsection{The management plane}

Every oversight mechanism deployed inside today's serving stacks shares a
structural property with the thing it oversees: it runs in the same
pipeline. Guard models classify the token stream in-line
\citep{llamaguard2023}; critique and reflection execute as part of the very
computation they are meant to check; representation-level monitors and
intervention methods, however promising, run in the process they observe
\citep{zou2024circuitbreakers}, and halt mechanisms reachable through the
model's own input channel inherit that channel's attack surface
\citep{killbench2025}. Execution-time safety kernels that treat the agent
runtime as untrusted and hold an out-of-reach stop \citep{safetykernel2026}
apply the same doctrine at the action layer; the management plane applies
it at the layer beneath, the resource contracts every action layer
depends on. The in-pipeline property is not an implementation accident,
it is a fate: whatever compromises, overloads, or wedges the inference
stack is positioned to compromise its oversight in the same stroke, and
the scenarios in which oversight matters most, a manipulated agent, a
saturated cluster, a defective rollout, are precisely the scenarios in
which the stack is least trustworthy about itself. Section 3's third
lesson states the doctrine networks derived from the same observation.

The discipline therefore places supervision in a management plane with
two independence properties, named exactly. The first is authority
independence: the plane holds pre-authorized rights over flows, and the
validation and audit of those rights occur outside the inference stack,
so no state of the model or the serving engine can expand, contract, or
forge them. The second is availability independence: the plane runs in a
separate failure domain, a distinct process and host at minimum, distinct
credentials and control network in the strong embodiment, with a
heartbeat between the planes; if the management plane becomes
unreachable, the serving layer enters a configured fail-safe, admitting
no new reservations while honoring existing guarantees, so that
supervision failure degrades the system toward caution rather than toward
the unsupervised open. We do not claim, and the design does not require,
signal independence: the plane's inputs, telemetry, ledger state, and
activation taps exported one-way from the data plane, originate in the
stack it watches, and a sufficiently compromised stack can starve or
distort them. What the compromised stack cannot do is revoke the
watcher's authority or take it offline with itself, and that is the
property the doctrine actually needs.

The plane's action vocabulary is small and maps onto the discipline's
objects. Halt revokes a reservation and stops a flow: the one override of
Section 4.2's guarantee, triggerable only by this authority and never by
a capacity condition. Steer forces a degradation action from Section
4.2's ladder onto a specific flow, a bounded intervention short of
termination. Escalate routes a flow's disposition to an external policy
engine or a human. Each action is pre-authorized at admission time,
logged against the ledger, and executable without the cooperation of the
flow, the model, or the scheduler that serves them.

The control plane as an enforcement point. It is worth stating plainly
what this machinery amounts to for the governance of deployed agents,
because the question of who can stop a misbehaving agent is today
answered mostly inside the model, through training, prompting, and
in-line classification, which is to say inside the thing being doubted.
The discipline gives an organization four primitives that require no
trust in the model at all. Admission is where policy binds: an
organization can require that flows touching designated systems present
an attested identity, a bounded budget, and an eligible provenance class
before any capacity is granted. The budget is containment: the worst
case of a rogue or looping agent is its remaining budget, a number fixed
at admission rather than discovered in an incident review, and Section
7's adversarial results measure exactly this property. The management
plane is the stop: a pre-authorized, audited, out-of-band kill that does
not route through the agent's own reasoning. And the reservation ledger
is the record: for every admitted flow, who authorized it, under what
identity and provenance, bounded to what, degraded or halted when and by
whom. Regulatory frameworks are beginning to require exactly this shape
of control: the EU AI Act's human-oversight article mandates an
intervention or stop capability for high-risk systems \citep{euaiact_art14},
and a reservation-holding control plane is where such a stop can be
implemented as a deterministic property of infrastructure rather than a
request to the model. These are the questions compliance functions
actually ask, and identity-layer controls can already revoke an agent's
credentials and tool permissions without the model's cooperation; the
serving layer adds the answer those controls lack, revocation of the
computation itself, since an agent with revoked credentials but live
compute continues reasoning toward its next attempt, while an agent with
a revoked reservation stops. The boundary of the claim is equally plain,
and Section 10 fixes it as a non-goal: this layer enforces decisions
about flows; it does not judge content. Detection, whether a flow's
outputs are harmful, deceptive, or noncompliant, is a model-layer problem
with its own literature and its own tools, and those tools gain rather
than lose value here: a verdict from any of them, or from a human,
becomes actionable through an authority that the misbehaving computation
cannot reach.

\subsection{Trust-keyed admission}

The discipline's final element keys the admission decision of Section 4.1
to provenance, the transfer of Section 3's fourth lesson. Every flow
carries a trust class assigned at admission from the provenance of the
content that drives it: an instruction typed by an authenticated
operator; a request from an authenticated service principal; a task
derived from tool output; a task derived from open-web content. The
classes order naturally, and the rule is structural: flows at or below a
configured threshold are ineligible for the reserved subset regardless of
their content, receive capped budgets, and sit lowest in the
class-differentiated degradation curves of Section 4.2, shedding first
under pressure.

The qualifier regardless of their content is the design. Content-level
defenses against prompt injection and instruction subversion are an
active field with real mechanisms: control-flow and data-flow separation
with capability policies over tool access \citep{camel2025}, information-flow
labels propagated through agent pipelines \citep{fides2025}, runtime policies
gating tool invocation on taint \citep{rtbas2025}, and trained instruction
hierarchies \citep{wallace2024}. These systems, with argument-level and
formally enforced descendants \citep{pact2026,provgates2026}, answer the
question of what a piece of content may cause an agent to do, and they
answer it downstream of admission, in the agent's execution; integrity
layers that verify authority at the moment of execution extend the same
family \citep{cxi2026}. Trust-thresholded admission itself has appeared for one
resource, entries into an agent's memory store \citep{memcurate2026}. The
discipline's object is different: serving-layer capacity and standing.
What provenance gates here is whether a flow may hold a guaranteed
reservation, at what budget scale, and where it stands in the degradation
order, fixed at admission before the flow's content is ever interpreted.
Composed with Section 4.3, the enforcement reading is direct: untrusted
provenance is served, at best-effort and under bound, but it cannot
occupy the infrastructure's promises.

Reservation survivability. A guarantee that dies with a serving instance
is a guarantee against everything except the most common failure. The
discipline extends the contract across unplanned instance failure by
replicating the reservation record itself, the budget counter and service
class, deliberately exclusive of cache contents and generation progress,
to a standby. On failure, the standby resumes the admitted flow without
re-entering admission and without enqueueing behind its own intake, and
it makes room for the inherited guarantee under load the same way the
discipline makes room for anything: by degrading its best-effort traffic
through the headroom-keyed ladder. The flow loses unpinned progress; it
does not lose its contract.
